\section{符号说明}

四个问题共用同一套符号。电量类变量以 kWh 为单位、本身已含时间因子；功率类输入以 kW 为单位；费用类量以元为单位。

\begin{center}
\small
\begin{tabular}{clc}
\toprule
符号 & 含义 & 单位 \\
\midrule
$t,\ \tau$ & 时段索引，每日 $t=1,\dots,144$；全年 $\tau=1,\dots,52560$ & --- \\
$d,\ i$ & 日期索引（全年 365 天）；日内 10 分钟时段索引 $i=1,\dots,144$ & --- \\
$m$ & 决策时刻，$m\in\{0,6,12,18\}$ & h \\
$\Delta t$ & 时段长度，$\Delta t=1/6$ & h \\
$L_t$ & 小区负载功率 & kW \\
$PV_t$ & 光伏发电功率；$PV^{fc}_t$ 为预报值，$PV^{act}_t$ 为实际值 & kW \\
$p_t$ & 电价（问题一至三取自附件 1，逐日重复） & 元/kWh \\
$p^{act}_{d,i}$ & 附件 4 的实际波动电价 & 元/kWh \\
$\hat p^{(m)}_{d,i}$ & 决策时刻 $m$ 信息集下的预测电价；$\kappa_m$ 为调整层的水平校正因子 & 元/kWh \\
$b_t$ & 计划购电量（0:00 制定） & kWh \\
$q_t$ & 调整后的最终购电量（替代量口径） & kWh \\
$q^{em}_t$ & 紧急购电量（缺口部分） & kWh \\
$c_t,\ q^{dis}_t$ & 储能充电量、放电量 & kWh \\
$s_t$ & 弃光量（光伏盈余中未被消纳的部分） & kWh \\
$E_t$ & 时段末储电量（状态变量） & kWh \\
$\eta$ & 充放电效率，$\eta_{ch}=\eta_{dis}=0.9$ & --- \\
$E_{\min},E_{\max}$ & 储电量运行区间，$1200$ 与 $10800$ & kWh \\
$E_{\text{init}}$ & 初始储电量（2025-01-01 0:00），6000 & kWh \\
$P_{\max}$ & 最大充放电功率，5000 & kW \\
$\alpha_{em}$ & 紧急购电电价倍数，$\alpha_{em}=5$ & --- \\
$\beta_{def},\beta_{over}$ & 偏差结算倍数，$\beta_{def}=0.5$，$\beta_{over}=1.5$ & --- \\
$C_{\text{plan}},C_{\text{adj}},C_{\text{em}}$ & 计划购电费、偏差结算费、紧急购电费 & 元 \\
$C_{\text{total}}$ & 总费用 $C_{\text{plan}}+C_{\text{adj}}+C_{\text{em}}$ & 元 \\
$C^{fc}$ & 同一组决策按决策层信念价计得的费用 & 元 \\
$N$ & 净负荷电量 $\sum_t(L_t-PV_t)\Delta t$ & kWh \\
$D_{\text{full}},D_{\text{req}}$ & 全年 365 天；交付期 2025-02-01 至 12-31 共 334 天 & --- \\
$T$ & 决策期时段总数（问题二为 52\,560） & --- \\
\bottomrule
\end{tabular}
\end{center}

需要特别说明两处容易混用、一旦混用就无法比对的量：其一，$b$ 与 $q$ 是两个不同的量，$b$ 是 0:00 制定的计划购电量，$q$ 是该时段\emph{最终}的购电量（替代量口径而非增量），问题三与链 \emph{4-3} 的交付表报的是 $q$；其二，“全期”指 365 天，“交付期”指 2025 年 2 月 1 日之后的 334 天，1 月作为储能预热期不计入交付——全文出现的每一个费用数字都标注其所属口径。
